\SourceLeaf{351}%
auxiliary sphere, which gives the triangle's total curvature, is to the entire surface of that sphere. In general, the excess of the angles of a polygon with \(n\) sides over \(2n-4\) right angles equals its total curvature if its sides are shortest lines.

At the end of the memoir the general investigations developed there are applied to triangles formed by shortest lines. We give only a couple of the main theorems here. Let \(a,b,c\) be the sides of such a triangle, regarded as quantities of first order, \(A,B,C\) the opposite angles, \(\alpha,\beta,\gamma\) the measures of curvature at the vertices, and \(\sigma\) the area of the triangle. Apart from quantities of fourth order, \(\tfrac13(\alpha+\beta+\gamma)\sigma\) is the excess of the sum \(A+B+C\) over two right angles. To the same accuracy, the angles of a plane straight-sided triangle with sides \(a,b,c\) are, in order,
\[
\begin{aligned}
&A-\tfrac1{12}(2\alpha+\beta+\gamma)\sigma\\
&B-\tfrac1{12}(\alpha+2\beta+\gamma)\sigma\\
&C-\tfrac1{12}(\alpha+\beta+2\gamma)\sigma
\end{aligned}
\]
It is immediately clear that the last theorem generalizes the familiar one first stated by \textsc{Legendre}. According to that theorem, apart from quantities of fourth order, one obtains the angles of the straight-sided triangle by reducing each angle of the spherical triangle by one third of the spherical excess. On a nonspherical surface one must therefore apply unequal reductions to the angles, and the inequality is generally a quantity of third order. If, however, the entire surface departs only slightly from a sphere, the inequality also involves a factor of the order of that departure. It is unquestionably important for higher geodesy to be able to calculate these inequalities and thereby be fully assured that they are imperceptible for all measurable triangles on Earth's surface. For example, in the largest triangle of the author's triangulation, whose largest side is nearly 15 geographical miles long and whose angle sum exceeds two right angles by nearly 15 seconds, the three reductions of the angles to %
